\section*{الفصل \ref{ch:b3:deformation-mechanisms} --- آليات التشوّه ونطاقات القصّ}

\begin{solution}{exo:b3:deformation-mechanisms:1}
الفورستيريت: $1573/2163 = 0.73$؛ الجليد: $263/273 = 0.96$. الجليد أقرب بكثير إلى
نقطة انصهاره ويزحف بسهولة أكبر: تنساب الأنهار الجليدية
على سطح الأرض.
\end{solution}

\begin{solution}{exo:b3:deformation-mechanisms:2}
$\gamma = \tan\ang{70} = 2.75$ و$\tan\ang{80} = 5.67$. بالقرب من \ang{90} يتغيّر الظلّ
بسرعة كبيرة مع الزاوية: خطأ بمقدار درجة واحدة عند $\psi =
\ang{84}$ يغيّر $\gamma$ بنحو الخُمس.
\end{solution}

\begin{solution}{exo:b3:deformation-mechanisms:3}
$\sigma = (10^{3.56}/10)^{1/1.26} = 363.1^{0.794} = \qty{108}{MPa}$.
\end{solution}

\begin{solution}{exo:b3:deformation-mechanisms:4}
$E/RT = 223000/(8.314\times773.15) = 34.69$، إذن $\dot\varepsilon =
\num{1.1e-4}\times20^4\times e^{-34.69} = \num{1.1e-4}\times\num{1.6e5}\times
\num{8.6e-16} = \qty{1.5e-14}{s^{-1}}$.
\end{solution}

\begin{solution}{exo:b3:deformation-mechanisms:5}
$\eta = \sigma/(2\dot\varepsilon)$: $\num{65.4e6}/\num{2e-14} =
\qty{3.3e21}{Pa.s}$ عند \qty{400}{\celsius} و$\num{18.0e6}/\num{2e-14} =
\qty{9.0e20}{Pa.s}$ عند \qty{500}{\celsius}: القشرة الوسطى عند هذه
درجات الحرارة تماثل تقريبًا \omterm{def:b3:mantle-convection:gia}{لزوجة الوشاح} الواقع تحت الأغطية الجليدية.
\end{solution}

\begin{solution}{exo:b3:deformation-mechanisms:6}
$n = \log(\num{1.6e-4}/\num{e-5})/\log2 = \log16/\log2 = 4$: \omterm{def:b3:deformation-mechanisms:creep}{الزحف الانخلاعي}،
الذي يتجاوز \omterm{def:b3:deformation-mechanisms:flow-law}{أسّ الإجهاد} فيه الواحد بكثير.
\end{solution}

\begin{solution}{exo:b3:deformation-mechanisms:7}
$\ln10 = (E/R)(1/1000 - 1/1100)$، إذن $E = 8.314\times2.3026/\num{9.091e-5} =
\qty{211}{kJ.mol^{-1}}$.
\end{solution}

\begin{solution}{exo:b3:deformation-mechanisms:8}
نابارو--هيرينغ ($d^{-2}$): أسرع بمقدار 100 مرة؛ كوبل ($d^{-3}$): أسرع بمقدار 1000 مرة.
على طول الحدود تسلك الذرات المنتشرة مسارًا مقطعه العرضي، $\delta
d$، يتناقص بتناقص الحبيبة، فوق المسافة الأقصر والحجم الأصغر
الواجب تحريكه.
\end{solution}

\begin{solution}{exo:b3:deformation-mechanisms:9}
$\sigma_b^{2.5} = \num{3.59e-6}\,d^{-3}/\num{6.54e-17}$: من أجل $d = 1000$،
$\sigma_b = 54.9^{0.4} = \qty{5.0}{MPa}$؛ ومن أجل $d = 100$، $\sigma_b =
\num{54900}^{0.4} = \qty{79}{MPa}$. الميل هو $-3/2.5 = -1.2$: تصغير الحبيبة
بعشرة أضعاف يرفع إجهاد الحدود بمقدار $10^{1.2} = 16$.
\end{solution}

\begin{solution}{exo:b3:deformation-mechanisms:10}
في المسارين كليهما يكون فرق التركيز المحرّك نفسه،
والمسافة هي $d$؛ والحجم الواجب نقله لكل وحدة انفعال هو $d^3$. توفّر الشبيكة
مقطعًا عرضيًا $\propto d^2$، فيتناسب المعدّل مع $d^2/(d\,d^3) =
d^{-2}$؛ وتوفّر الحدود $\propto\delta d$، فيكون $\delta d/(d\,d^3) = \delta
d^{-3}$. نسبة معدّلي الحدود إلى الشبيكة هي $(B_B\delta D_B)/(B_LD_Ld)$،
أكبر من 1 عندما $d < (B_B/B_L)\,\delta D_B/D_L$.
\end{solution}

\begin{solution}{exo:b3:deformation-mechanisms:11}
$\eta \propto \dot\varepsilon^{(1-n)/n} = \dot\varepsilon^{-2.5/3.5}$: بمعدّل أسرع بعشر
مرات تُقسم اللزوجة على $10^{0.714} = 5.2$. النطاق الذي يتّفق أن
يتشوّه بسرعة أكبر يصبح أضعف، فيحمل النصيب الأكبر من التشوّه: يتجمّع
الانفعال فيه.
\end{solution}

\begin{solution}{exo:b3:deformation-mechanisms:12}
قبل: $\dot\varepsilon_{\mathrm{dis}} = \num{6.54e-17}\times10^{3.5} =
\qty{2.07e-13}{s^{-1}}$ و$\dot\varepsilon_{\mathrm{dif}} = \num{3.59e-6}
\times10/3000^3 = \qty{1.3e-15}{s^{-1}}$: \qty{2.08e-13}{s^{-1}}، \omterm{def:b3:deformation-mechanisms:creep}{الزحف
الانخلاعي}. بعد: $\dot\varepsilon_{\mathrm{dif}} = \num{3.59e-6}\times10/30^3 =
\qty{1.33e-9}{s^{-1}}$، \omterm{def:b3:deformation-mechanisms:creep}{الزحف الانتشاري}، أسرع بمقدار \num{6400} مرة عند
الإجهاد نفسه: ليّنت إعادة التبلور النطاق تليينًا كبيرًا.
\end{solution}

\begin{solution}{pb:b3:deformation-mechanisms:1}
\textbf{1.} يدور القاطع من \ang{90} إلى \ang{6} بالنسبة إلى النطاق: $\psi =
\ang{84}$ و$\gamma = \tan\ang{84} = 9.51$.

\textbf{2.} $\gamma w = 9.51\times400 = \qty{3804}{m}$، أي نحو \qty{3.8}{km}.

\textbf{3.} $S_1 = \sqrt{1 + 9.51^2/4} + 9.51/2 = 4.859 + 4.755 = 9.614$,
$S_2 = 1/S_1 = 0.104$، و$S_1/S_2 = 92.4$.

\textbf{4.} $\tan2\theta' = 2/9.51 = 0.210$، $\theta' = \ang{5.9}$: يدور
التورّق، القريب من المحور الطويل، يقترب أكثر فأكثر من النطاق كلما $\gamma$
تزايدت.

\textbf{5.} عند تتبّعه نحو الشمال يُحمل القاطع شرقًا داخل النطاق،
ويغادره أبعد نحو الشرق: انتقل الجدار الشمالي شرقًا بالنسبة إلى الجنوبي.
عند النظر من الجدار الجنوبي تحرّك الجانب البعيد نحو اليمين: النطاق
أيمن.

\textbf{6.} في المستوي الذي يحوي تخطّط المطّ، أفقيًا هنا
لنطاق انزياحي، وعموديًا على التورّق الرأسي: مقطع
أفقي. يميل تورّق S بعكس اتجاه عقارب الساعة عن النطاق، باتجاه
شرق-شمال-شرق، مثل المحور الطويل للقطع الناقص للانفعال.

\textbf{7.} $\sigma = (10^{3.56}/20)^{1/1.26} = 181.5^{0.794} = \qty{62}{MPa}$.

\textbf{8.} $242.1^{0.794} = \qty{78}{MPa}$ من أجل \qty{15}{\micro m}
و$145.2^{0.794} = \qty{52}{MPa}$ من أجل \qty{25}{\micro m}.

\textbf{9.} نعم: امتدّت التجارب من 34 إلى \qty{268}{MPa}.

\textbf{10.} تشوّهت المنطقة عند \qty{450}{\celsius}، أدنى بكثير من 800 إلى
\qty{1100}{\celsius} في التجارب؛ علاقةٌ لا تعتمد على
درجة الحرارة يمكن تطبيقها عليها، وإن كانت ما تزال
استقراءً.

\textbf{11.} $D = 3631/100^{1.26} = 3631/331 = \qty{11}{\micro m}$.

\textbf{12.} $E/RT = 223000/(8.314\times723.15) = 37.09$ و$e^{-37.09} =
\num{7.79e-17}$.

\textbf{13.} $\dot\varepsilon = \num{1.1e-4}\times62^4\times\num{7.79e-17} =
\num{1.1e-4}\times\num{1.478e7}\times\num{7.79e-17} = \qty{1.27e-13}{s^{-1}}$.

\textbf{14.} $\dot\gamma = 2\dot\varepsilon = \qty{2.53e-13}{s^{-1}}$.

\textbf{15.} $\eta = \num{62e6}/(2\times\num{1.27e-13}) = \qty{2.4e20}{Pa.s}$,
قريبةً من \qty{3e20}{Pa.s} للوشاح العلوي تحت بحيرة آردو: كانت المنطقة
ضعيفةً بقدر الوشاح.

\textbf{16.} يتغيّر $e^{-E/RT}$ بالعامل 0.064 عند \qty{400}{\celsius}
(أبطأ ستّ عشرة مرّة)، و11.0 عند \qty{500}{\celsius}.

\textbf{17.} $(78/62)^4 = 2.5$ مرّة أسرع.

\textbf{18.} $t = 9.51/\num{2.53e-13} = \qty{3.75e13}{s}$.

\textbf{19.} $\num{3.75e13}/\num{3.156e10} = \qty{1190}{kyr}$.

\textbf{20.} $\qty{3804}{m}/\qty{1.19e6}{yr} = \qty{3.2}{mm/yr}$.

\textbf{21.} أبطأ بعشر مرّات أو أكثر من حدود الصفائح الرئيسية: كانت المنطقة
فرعًا ثانويًا، أو واحدًا من عدّة فروع تشترك في حركة حدٍّ.

\textbf{22.} $\dot\gamma = 0.02/(\num{3.156e7}\times400) = \qty{1.58e-12}{s^{-1}}$،
6.2 مرّة أسرع؛ يرتفع الإجهاد بمقدار $6.2^{1/4} = 1.58$، إلى \qty{98}{MPa}،
وستكون الحبيبات $3631/98^{1.26} = \qty{11}{\micro m}$.

\textbf{23.} يتناسب الزمن مع $\sigma^{-4}$: يعطي مدى حجم الحبيبات وحده
$(62/78)^4 = 0.40$ إلى $(62/52)^4 = 2.0$ مرّةً من التقدير،
وعدم اليقين في الأس $1.26 \pm 0.13$ وفي ثوابت \omterm{def:b3:deformation-mechanisms:flow-law}{قانون الانسياب}
يوسّعه أكثر: النتيجة رتبة قدر، لا تاريخًا.

\textbf{24.} استغرق \omterm{def:b3:deformation-mechanisms:shear-zone}{نطاق القصّ} في تلال فاردن نحو
\textbf{\qty[text-series-to-math=true]{1190}{kyr}} لتراكم انفعاله.
\end{solution}
